\begin{abstract}
% skip one line to make the abstract start with indent
For decades, the evolution of operating systems has been shaped by a proven conceptual model that emerged in the mid-20th century with the advent of UNIX. Today, however, this model is increasingly revealing its fundamental limitations. Attempts to circumvent these limitations as early as the late 20th century led, among other things, to the emergence of the concept of operating systems with persistent memory. Despite the merits of this paradigm, research was effectively halted due to hardware limitations. Since then, computer performance has increased significantly, giving impetus to new work in this direction. A striking illustration of this is the Phantom OS. However, writing and implementing a completely new operating system is a lengthy, labor-intensive process that entails a large number of errors. Mitigating these consequences is aided by the use of proven mechanisms, which is facilitated by porting to Genode. In this paper, we examine the process of adapting PhantomOS components for Genode. The issue of secure snapshot storage remains relevant. To address this, file\_vault was adapted to perform block-level encryption and provide a file system session. This led to the creation of se\_vault. However, the downside of this security measure, which limits its usability, is a significant performance degradation—approximately 24 times slower—due to the use of the experimental Tresor library. These results open up avenues for further research in the field of persistent operating system security and may contribute to their wider adoption.

\end{abstract}